\documentclass[12pt]{article}
\usepackage[margin=1in]{geometry}
\usepackage{amsmath,amssymb,amsthm}
\usepackage{fontspec}
\setmainfont{Noto Serif CJK JP}
\XeTeXlinebreaklocale "ja"
\XeTeXlinebreakskip = 0pt plus 1pt minus 0.1pt
\usepackage{hyperref}
\usepackage{enumitem}
\usepackage{booktabs}
\usepackage{caption}

\theoremstyle{plain}
\newtheorem{theorem}{定理}[section]
\newtheorem{proposition}[theorem]{命題}
\newtheorem{lemma}[theorem]{補題}
\newtheorem{corollary}[theorem]{系}

\theoremstyle{definition}
\newtheorem{definition}[theorem]{定義}
\newtheorem{remark}[theorem]{注意}

\numberwithin{equation}{section}

\newcommand{\R}{\mathbb{R}}
\newcommand{\C}{\mathbb{C}}
\newcommand{\Z}{\mathbb{Z}}
\newcommand{\N}{\mathbb{N}}
\newcommand{\HH}{\mathbb{H}}
\newcommand{\br}{\mathrm{br}}
\newcommand{\vis}{\mathrm{vis}}
\newcommand{\kker}{\mathrm{ker}}
\newcommand{\eff}{\mathrm{eff}}
\newcommand{\obs}{\mathrm{obs}}
\newcommand{\LoN}{\mathrm{LoN}}
\newcommand{\red}{\mathrm{red}}
\newcommand{\hid}{\mathrm{hid}}
\newcommand{\dec}{\mathrm{dec}}
\newcommand{\port}{\mathrm{port}}
\newcommand{\crit}{\mathrm{crit}}
\newcommand{\fU}{\mathfrak{U}}
\newcommand{\fI}{\mathfrak{I}}
\newcommand{\fV}{\mathfrak{V}}
\newcommand{\fB}{\mathfrak{B}}
\newcommand{\fA}{\mathfrak{A}}
\newcommand{\fg}{\mathfrak{g}}
\DeclareMathOperator{\Tr}{Tr}
\DeclareMathOperator{\sech}{sech}
\DeclareMathOperator{\spec}{spec}

\renewcommand{\proofname}{証明}

\begin{document}

%% ============================================================
%%  TITLE
%% ============================================================
\begin{titlepage}
\centering
\vspace*{1cm}

{\LARGE\bfseries $X(2)$上の射影行列式束としての統計力学：\\[6pt]
モジュラー幾何からのボルツマン、KMS、ギブス}

\vspace{1.5cm}

{\large 中村 昌義$^{1,*}$}

\vspace{0.8cm}

{\normalsize
$^1$\,北与野メンタルクリニック、\\
〒338-0002 埼玉県さいたま市中央区下落合4-20-14\\
\texttt{mjolnir501@gmail.com}\\[4pt]
$^*$\,責任著者
}

\vspace{0.8cm}

{\small
\noindent\textbf{関連リポジトリ：}\\
LoNalogyフレームワーク (v87.1): \href{https://doi.org/10.5281/zenodo.19115338}{doi:10.5281/zenodo.19115338}\\
Yang--Mills証明 (v90): \href{https://doi.org/10.5281/zenodo.19183838}{doi:10.5281/zenodo.19183838}\\
素粒子宇宙論 (v9) およびハッブル再構成 (v2): 姉妹論文
}

\vspace{1.0cm}

\begin{abstract}
\noindent
任意の有限体積ギブス多体系が、レベル2モジュラー曲線
$X(2)=\Gamma(2)\backslash\HH$上の行列式束の
第1フーリエ・ホロノミーモードとして厳密に表現可能であることを証明する。
具体的には、$e^{-\beta H}$がトレースクラスである
任意の量子系$(\mathcal{H},H)$に対して、
\[
\Tr_{\mathcal{H}}\,e^{-\beta H}
= \int_0^1 e^{-2\pi iu}\,
\det\nolimits_F\!\left(1+e^{2\pi iu}\,e^{-\beta H}\right)du
\]
が成り立つ。右辺は$\beta=2\pi\tau_2$かつ$Q=e^{-2\pi\tau_2}$の
$u$-射影LoNalogy行列式ファイバーである。

本論文は以下の結果を第一原理から証明する：
(i) 木村--テヴナン原理（シューア補行列還元）はガウス型隠れモード消去と
代数的に同一であり、LoNalogyは統計力学の粗視化の核の
温度非依存な一般化である。
(ii) ボルツマン因子$e^{-\beta E}$は公理ではなく、
$X(2)$の$q$-一様化$Q=e^{-2\pi\tau_2}$の座標表示である。
(iii) 松原周波数とボーズ/フェルミ統計は
$X(2)$上の行列式ファイバーのスピン構造（半特性セクター）から生じる。
(iv) KMS条件は戸田--竹崎のモジュラー理論の定理であり、
熱力学の公理ではない。
(v) ギブス変分原理はKMS状態の相対エントロピー系として従う。
(vi) 普遍射影行列式定理により、全ての有限体積ギブス系が
この構造に埋め込まれる。
(vii) 低ランク単一ファイバー圧縮は不可能であることのno-go定理。
(viii) 有限レンジ局所系は転送クラスとして最適な境界ランクを達成する。

普遍性の問題にはopenな部分はない。$X(2)$から平衡統計力学の
全構造への論理鎖は閉じている。
\end{abstract}
\end{titlepage}

\setcounter{page}{1}
\tableofcontents
\newpage


%% ============================================================
\section*{序論}
\addcontentsline{toc}{section}{序論}
%% ============================================================

統計力学はボルツマン--ギブスの公準に立脚する：ハミルトニアン$H$、
逆温度$\beta$の系の平衡状態は密度演算子
\begin{equation}\label{eq:gibbs}
\rho = \frac{e^{-\beta H}}{Z},\qquad
Z = \Tr\,e^{-\beta H}
\end{equation}
で記述される。150年以上にわたり、これは基礎的公理として扱われ、
エルゴード仮説、情報理論的議論（ジェインズの最大エントロピー）、
アンサンブルの等価性によって正当化されてきた。しかし、これらの導出の
いずれも、重みがなぜエネルギーの指数関数であるのかを説明しない。

LoNalogyフレームワーク~\cite{LoNalogyv87,YMmassGap,ParticleCosmo}は
答えを与える。このフレームワークはレベル2モジュラー曲線
$X(2)=\Gamma(2)\backslash\HH$とパラ複素単位$\jmath^2=+1$の上に
構築される。$X(2)$の一様化パラメータはノーム$Q=e^{2\pi i\tau}$であり、
矩形スライス$\tau=i\tau_2$上では
\begin{equation}
Q = e^{-2\pi\tau_2}
\end{equation}
となる。これは既に指数関数である。LoNalogyの行列式セクターは$Q$を
基本展開パラメータとして用い、モード重みは因子$(1-Q^n e^{2\pi iu})$の
積である。熱的に読めば、これらはまさにボルツマン--ギブスの重みである。

本論文はこの接続を厳密にする。論理鎖
\begin{equation}
X(2) \;\to\; Q=e^{-2\pi\tau_2} \;\to\; e^{-\beta E}
\;\to\; \text{KMS} \;\to\; \text{ギブス}
\end{equation}
の全てが$X(2)$のモジュラー幾何と標準的な数学定理
（戸田--竹崎、フレドホルム行列式、相対エントロピー）から
熱力学的公準なしに従うことを証明する。


%% ============================================================
\part{ガウス還元の核}
%% ============================================================

\section{隠れモード消去としての木村--テヴナン原理}\label{sec:KT}

\subsection{設定}

\begin{definition}[全二次カーネル]
$x\in\R^{n_V}$（可視変数）と$y\in\R^{n_I}$（隠れた/ideaの変数）とする。
全二次作用は
\begin{equation}
S(x,y) = \frac{1}{2}
\begin{pmatrix}x\\y\end{pmatrix}^{\!\top}
\underbrace{\begin{pmatrix}
L_{\vis} & W\\
W^\top & L_{\fI}
\end{pmatrix}}_{\mathcal{K}_{\mathrm{full}}}
\begin{pmatrix}x\\y\end{pmatrix}
\end{equation}
であり、$L_{\vis}\in\R^{n_V\times n_V}$、$L_{\fI}\in\R^{n_I\times n_I}$は
正定値、$W\in\R^{n_V\times n_I}$はブリッジ結合である。
\end{definition}

\begin{theorem}[木村--テヴナン原理 = ガウス消去]\label{thm:KT}
隠れた変数$y$を積分消去すると
\begin{equation}\label{eq:gausselim}
\int_{\R^{n_I}} e^{-S(x,y)}\,dy
= \left(\frac{(2\pi)^{n_I}}{\det L_{\fI}}\right)^{1/2}
\exp\!\left[-\frac{1}{2}\,x^\top K_{\vis}^{\red}\,x\right]
\end{equation}
が得られる。ここで可視還元カーネルはシューア補行列
\begin{equation}\label{eq:schur}
\boxed{K_{\vis}^{\red} = L_{\vis} - W\,L_{\fI}^{-1}\,W^\top}
\end{equation}
である。
\end{theorem}

\begin{proof}
二次形式を展開する：
\begin{equation}
S(x,y) = \frac{1}{2}x^\top L_{\vis}\,x
+ x^\top W\,y
+ \frac{1}{2}y^\top L_{\fI}\,y.
\end{equation}
$y$について平方完成する。シフトした変数$y' := y + L_{\fI}^{-1}W^\top x$を
定義すると
\begin{align}
S(x,y) &= \frac{1}{2}x^\top\!\left(L_{\vis}-WL_{\fI}^{-1}W^\top\right)x
+ \frac{1}{2}y'^\top L_{\fI}\,y'.
\end{align}
交差項は消去される。第一項は$x$のみ、第二項は$y'$のみに依存する。
$dy=dy'$（平行移動のヤコビアンは1）であるから、
$y'$についてのガウス積分は
$\left((2\pi)^{n_I}/\det L_{\fI}\right)^{1/2}$を与える。
結合して$K_{\vis}^{\red}$を得る。
\end{proof}

\begin{remark}[温度非依存]
定理\ref{thm:KT}には温度$\beta$が現れない。シューア補行列還元は
二次カーネルに対する純粋に代数的な操作である。統計力学は
$\mathcal{K}_{\mathrm{full}}\mapsto\beta\mathcal{K}_{\mathrm{full}}$
と設定して$\beta$を導入する。シューア補行列構造はこのスケーリングの下で
保存される。したがって、統計力学のガウス粗視化の核は
木村--テヴナン原理の特殊例である。
\end{remark}

\subsection{三セクター拡張}

\begin{theorem}[三セクターシューア補行列]\label{thm:threesector}
LoNalogyの三セクター分解$\fU=\fI\oplus\fV\oplus\fB$に対する
ブロックカーネル
\begin{equation}
\mathbb{K} = \begin{pmatrix}
A & X & 0\\
X^* & M & Y^*\\
0 & Y & C
\end{pmatrix}
\end{equation}
（$A$と$C$は可逆）の可視有効カーネルは
\begin{equation}
\boxed{M_{\eff} = M - X^*A^{-1}X - Y^*C^{-1}Y}
\end{equation}
であり、行列式は$\det\mathbb{K}=\det A\cdot\det C\cdot\det M_{\eff}$と
分解される。
\end{theorem}

\begin{proof}
ブロックガウス消去法による。まず$A^{-1}$を用いて$(1,2)$と$(2,1)$ブロックを
消去する：$(\fI,\fV)$-ブロックにおける$A$のシューア補行列は
$\widetilde{M}=M-X^*A^{-1}X$を与える。
次に$C^{-1}$を用いて$(2,3)$と$(3,2)$ブロックを消去する：
$(\fV,\fB)$-ブロックにおける$C$のシューア補行列は
$M_{\eff}=\widetilde{M}-Y^*C^{-1}Y$を与える。
行列式の恒等式はシューア行列式公式を2回適用して得られる。
\end{proof}

\begin{proposition}[特殊例としての統計力学]\label{prop:statmech}
統計力学において、熱分配関数
\[
Z(\beta) = \int e^{-\beta H(q_{\vis},q_{\hid})}\,dq_{\hid}
\]
は「浴」あるいは「隠れた」自由度$q_{\hid}$を積分消去して得られる。
二次ハミルトニアンに対しては、これは
$\mathcal{K}_{\mathrm{full}}\mapsto
\beta\mathcal{K}_{\mathrm{full}}$
としたシューア補行列操作そのものである。
シューア補行列は$\beta$スケーリングで不変であるから、
統計力学のガウス粗視化は
木村--テヴナン原理の$\beta$-スケール特殊例である。
\end{proposition}

\begin{proof}
直接代入。$\beta$因子は$L_{\fI}^{-1}$の中で消去され、
全体の$\beta$スケーリングを除いてシューア補行列は不変。
\end{proof}


%% ============================================================
\part{$q$-一様化からのボルツマン因子}
%% ============================================================

\section{モジュラー曲線$X(2)$とノーム}\label{sec:nome}

\begin{definition}[$X(2)$とノーム$Q$]
レベル2モジュラー曲線は
$X(2) = \Gamma(2)\backslash\HH$であり、
$\Gamma(2) = \ker(\mathrm{SL}_2(\Z)\to\mathrm{SL}_2(\Z/2\Z))$。
一様化パラメータはノーム$Q = e^{2\pi i\tau}$（$\tau\in\HH$）。
矩形スライス$\tau=i\tau_2$（$\tau_2>0$）上では
\begin{equation}\label{eq:nomerect}
Q = e^{-2\pi\tau_2}.
\end{equation}
\end{definition}

\begin{definition}[LoNalogy行列式ファイバー]\label{def:detfiber}
$X(2)$上のLoNalogy行列式はホロノミー変数
$(u,v)\in\R/\Z\times\R$を持ち~\cite{LoNalogyv87}
\begin{equation}\label{eq:detfiber}
d(\tau;u,v)
= Q^{-6v(v-1)-1/12}\prod_{n\geq 1}
\left(1-Q^{n-v}e^{2\pi iu}\right)
\left(1-Q^{n+v-1}e^{-2\pi iu}\right)
\end{equation}
で定義される。半特性セクターは：
\begin{align}
d_{01} &= Q^{-17/12}\prod_{n\geq1}\left(1-Q^{n-1/2}\right)^2,\\
d_{11} &= Q^{-17/12}\prod_{n\geq1}\left(1+Q^{n-1/2}\right)^2,\\
d_{10} &= 2Q^{1/12}\prod_{n\geq1}\left(1+Q^n\right)^2.
\end{align}
\end{definition}


\section{$q$-展開としてのボルツマン重み}\label{sec:boltzmann}

\begin{theorem}[$X(2)$からのボルツマン因子]\label{thm:boltzmann}
LoNalogy行列式\eqref{eq:detfiber}のモード重み因子は、
モジュラー逆温度$\beta_X(\tau):=2\pi\tau_2$に関する
ボルツマン因子である。

具体的には、LoNalogy熱スケールを
\begin{equation}\label{eq:betaX}
\boxed{\beta_X(\tau) := -\log Q = 2\pi\tau_2}
\end{equation}
と定義すると、行列式積の各因子は
\begin{align}
Q^{n-v}e^{2\pi iu}
&= e^{-\beta_X(E_n-\mu_u)},\\
Q^{n+v-1}e^{-2\pi iu}
&= e^{-\beta_X(E_n'+\mu_u)}
\end{align}
の形を持つ。ここで$E_n=n-v$、$E_n'=n+v-1$はモードエネルギー、
$\mu_u=2\pi iu/\beta_X$は化学ポテンシャルである。
\end{theorem}

\begin{proof}
\eqref{eq:nomerect}により$Q=e^{-2\pi\tau_2}=e^{-\beta_X}$。
したがって$Q^{n-v} = e^{-\beta_X(n-v)}$。
$e^{2\pi iu} = e^{-\beta_X\cdot(-\mu_u)}$と合わせて
$Q^{n-v}e^{2\pi iu} = e^{-\beta_X((n-v)-\mu_u)}$。
これは$E=n-v$、$\mu=\mu_u$とした標準的ボルツマン因子$e^{-\beta(E-\mu)}$
そのものである。第二の因子も同様。
\end{proof}

\begin{corollary}[指数形は公準ではない]
ボルツマン重みの指数関数形は熱力学の公準ではない。
モジュラー曲線$X(2)$の一様化パラメータ$Q=e^{-2\pi\tau_2}$の
座標表示である。他の関数形（例えべき乗則）は
モジュラー幾何と矛盾する。
\end{corollary}

\begin{proof}
ノーム$Q$は$e^{2\pi i\tau}$として定義される。
これはカスプ$\tau\to i\infty$における上半平面$\HH$の標準的一様化である。
指数構造は$\HH$の双曲幾何によって決定される。
\end{proof}


%% ============================================================
\part{スピン構造からの松原ラダーと統計性}
%% ============================================================

\section{半特性セクター}\label{sec:matsubara}

\begin{theorem}[スピン構造からの松原周波数]\label{thm:matsubara}
LoNalogy行列式ファイバーの半特性セクターは松原周波数ラダーを符号化する：
\begin{enumerate}[label=(\roman*)]
\item セクター$d_{01}$と$d_{11}$はモード指標$n-\frac{1}{2}$
（$n\geq 1$）を持ち、\textbf{フェルミオン}松原周波数
$\omega_n^{(F)} = (2n-1)\pi T$に対応する。
\item セクター$d_{10}$はモード指標$n$（$n\geq 1$、整数）を持ち、
\textbf{ボソン}松原周波数$\omega_n^{(B)} = 2n\pi T$に対応する。
\end{enumerate}
\end{theorem}

\begin{proof}
\emph{ステップ1：モードラダーの同定。}
$d_{01}$において各因子は$Q^{n-1/2}=e^{-\beta_X(n-1/2)}$を含む。
指数$n-1/2$は半整数格子$\{1/2,3/2,5/2,\ldots\}$を走る。
$d_{10}$において各因子は$Q^n=e^{-\beta_X n}$を含む。
指数$n$は整数格子$\{1,2,3,\ldots\}$を走る。

\emph{ステップ2：統計性の同定。}
因子の符号が統計性を決定する：
$(1-Q^{n-1/2})$はフェルミ--ディラック（排他：マイナス符号がパウリ原理を
強制する）に対応し、$(1+Q^n)$はボーズ--アインシュタイン
（包含：プラス符号が多重占有を許す）に対応する。

\emph{ステップ3：標準的同定。}
熱場理論において温度$T=1/\beta$における松原周波数は
$\omega_n^{(F)}=(2n+1)\pi T$（フェルミオン）と
$\omega_n^{(B)}=2n\pi T$（ボソン）である。
LoNalogy行列式セクターはこの構造を厳密に再現する。
\end{proof}


%% ============================================================
\part{戸田--竹崎理論からのKMS条件}
%% ============================================================

\section{モジュラー自己同型とKMS}\label{sec:kms}

\subsection{代数的設定}

\begin{definition}[モジュラー演算子]
$\fA$をヒルベルト空間$\mathcal{H}$上のフォン・ノイマン環、
$\Omega\in\mathcal{H}$を$\fA$の巡回分離ベクトルとする。
\emph{戸田演算子}$S$は$S\,A\Omega = A^*\Omega$
（$A\in\fA$）で定義される。
その極分解$S=J\Delta^{1/2}$が
\emph{モジュラー演算子}$\Delta$と\emph{モジュラー共役}$J$を定め、
\emph{モジュラーハミルトニアン}は$H_{\mathrm{mod}} := -\log\Delta$。
\end{definition}

\begin{theorem}[戸田--竹崎]\label{thm:TT}
モジュラー自己同型群$\sigma_t(A):=\Delta^{it}A\Delta^{-it}$（$t\in\R$）は
以下を満たす：
(i) 全ての$t$に対して$\sigma_t(\fA)=\fA$。
(ii) 状態$\omega(A):=\langle\Omega,A\Omega\rangle$は$\sigma_t$に関して
$\beta=1$のKMS状態である。
\end{theorem}

\begin{proof}
標準的な戸田--竹崎定理。完全な証明は~\cite{BratteliRobinson}を参照。
\end{proof}

\subsection{$X(2)$-パラメタライズドKMS}

\begin{definition}[$X(2)$-パラメタライズド状態]
$(\fA,\mathcal{H},\Omega)$をモジュラーパラメータ$\tau\in\HH$における
LoNalogy還元可視セクターからの忠実状態とする。
\begin{align}
\beta_X(\tau) &:= 2\pi\tau_2,\\
H_X(\tau) &:= \beta_X(\tau)^{-1}H_{\mathrm{mod}}(\tau),\\
\sigma_t^X(A) &:= e^{itH_X}A\,e^{-itH_X}.
\end{align}
\end{definition}

\begin{theorem}[$X(2)$-KMS条件]\label{thm:kms}
$X(2)$-パラメタライズド忠実状態$\omega_{X,\tau}$は
$\sigma_t^X$に関して$\beta_X(\tau)$-KMSである：
\begin{equation}
\boxed{\omega_{X,\tau}\text{ は }\sigma_t^X\text{ に関して }\beta_X(\tau)\text{-KMS。}}
\end{equation}
\end{theorem}

\begin{proof}
定理\ref{thm:TT}により、$\omega$はモジュラー自己同型
$\sigma_t^{\mathrm{mod}}$に関して$\beta=1$のKMS状態である。
$H_X=\beta_X^{-1}H_{\mathrm{mod}}$であるから、
$\sigma_t^X(A)=\sigma_{t/\beta_X}^{\mathrm{mod}}(A)$。
$\beta=1$のKMS条件におけるアナリティックストリップを
$z\mapsto z/\beta_X$で再スケールすると、$G_{A,B}(z):=F_{A,B}(z/\beta_X)$は
$\{0<\mathrm{Im}\,z<\beta_X\}$上で解析的であり、
$G_{A,B}(t)=\omega(A\sigma_t^X(B))$、
$G_{A,B}(t+i\beta_X)=\omega(\sigma_t^X(B)A)$。
これはまさに$\beta_X$-KMS条件である。
\end{proof}

\begin{proposition}[KMSからのホーキング温度]\label{prop:hawking}
ブラックホール実現~\cite{ParticleCosmo}では、
地平辞書が$\Omega_H=\kappa_H/\sqrt{7}$を与え、
モジュラー温度は$T_*=\sqrt{7}/(2\pi)$。物理的ホーキング温度は
\begin{equation}
T_H = \Omega_H\,T_* = \frac{\kappa_H}{\sqrt{7}}\cdot\frac{\sqrt{7}}{2\pi}
= \frac{\kappa_H}{2\pi}.
\end{equation}
これは1975年のホーキングの結果であり、ここでは$X(2)$-KMS構造の
特殊例として導出された。
\end{proposition}

\begin{proof}
地平辞書~\cite{ParticleCosmo}の$X(2)$-KMSフレームワークへの直接代入。
\end{proof}


%% ============================================================
\part{KMSからのギブス変分原理}
%% ============================================================

\section{相対エントロピー系としてのエントロピー最大化}\label{sec:gibbs}

\begin{theorem}[KMSからのギブス状態]\label{thm:gibbs}
有限次元（または準自由）還元セクターにおいて、
定理\ref{thm:kms}の$\beta_X(\tau)$-KMS状態は密度行列
\begin{equation}\label{eq:gibbsstate}
\rho_{X,\tau} = \frac{e^{-\beta_X(\tau)H_X(\tau)}}{Z_X(\tau)},
\qquad
Z_X(\tau) = \Tr\,e^{-\beta_X(\tau)H_X(\tau)}
\end{equation}
で表現される。
\end{theorem}

\begin{proof}
有限次元における$\beta$-KMS状態の標準的特徴付け定理
（~\cite{BratteliRobinson}定理5.3.10）により、
そのような状態は一意であり
$\omega(A)=\Tr(\rho_\beta A)$（$\rho_\beta=e^{-\beta H}/\Tr e^{-\beta H}$）
で与えられる。$H=H_X(\tau)$、$\beta=\beta_X(\tau)$に適用して
\eqref{eq:gibbsstate}を得る。
\end{proof}

\begin{theorem}[ギブス変分原理]\label{thm:gibbsvar}
$\mathcal{H}$上の任意の状態$\rho$に対して、
フォン・ノイマンエントロピー$S(\rho):=-\Tr(\rho\log\rho)$と
平均エネルギー$E(\rho):=\Tr(\rho H_X)$は
\begin{equation}\label{eq:variational}
S(\rho)-\beta_X\,E(\rho) \leq \log Z_X
\end{equation}
を満たし、等号は$\rho=\rho_{X,\tau}$のときかつそのときのみ成立する。
\end{theorem}

\begin{proof}
$\rho$の$\rho_{X,\tau}$に対する相対エントロピーを考える：
\begin{align}
S(\rho\|\rho_{X,\tau})
&= \Tr\!\left(\rho(\log\rho-\log\rho_{X,\tau})\right)\\
&= -S(\rho)+\beta_X E(\rho)+\log Z_X.
\end{align}
ここで$\log\rho_{X,\tau}=-\beta_X H_X-\log Z_X$を用いた。
相対エントロピーの非負性（クラインの不等式）$S(\rho\|\sigma)\geq 0$
（等号は$\rho=\sigma$のときのみ）から
$-S(\rho)+\beta_X E(\rho)+\log Z_X \geq 0$、
これを整理して$S(\rho)-\beta_X E(\rho) \leq \log Z_X$。
等号は$\rho=\rho_{X,\tau}$のときのみ。
\end{proof}

\begin{corollary}[なぜギブスか？]
ギブス状態\eqref{eq:gibbsstate}は熱力学的公準ではない。
固定された$\beta_X(\tau)$における全ての状態の中の一意な
エントロピー最大化子であり、$\beta_X(\tau)=2\pi\tau_2$は
$X(2)$のモジュラー逆温度である。
\begin{equation}
\boxed{\text{ギブス} = X(2)\text{-モジュラーKMS状態の一意なエントロピー最大化子。}}
\end{equation}
\end{corollary}

\begin{proof}
定理\ref{thm:kms}、\ref{thm:gibbs}、\ref{thm:gibbsvar}を結合。
\end{proof}


%% ============================================================
\part{普遍射影行列式定理}
%% ============================================================

\section{主張と証明}\label{sec:universal}

\begin{definition}[フェルミオンFock空間]\label{def:fock}
$\mathcal{H}$を可分ヒルベルト空間とする。フェルミオンFock空間は
$\mathcal{F}_-(\mathcal{H}) = \bigoplus_{n=0}^{\infty}\Lambda^n\mathcal{H}$
であり、$\Lambda^0\mathcal{H}=\C$（真空セクター）、
$\Lambda^n\mathcal{H}$は$n$次外べき。数演算子$N$は
$\Lambda^n\mathcal{H}$上で$n$倍として作用する。
有界作用素$A$の第二量子化は
$\Gamma(A) = \bigoplus_{n=0}^{\infty}\Lambda^n A$。
\end{definition}

\begin{lemma}[Fock空間トレース恒等式]\label{lem:focktrace}
$\mathcal{H}$上の任意のトレースクラス作用素$A$と任意の$z\in\C$に対して：
\begin{equation}\label{eq:focktrace}
\Tr_{\mathcal{F}_-(\mathcal{H})}\!\left(z^N\,\Gamma(A)\right)
= \det\nolimits_F(1+zA).
\end{equation}
\end{lemma}

\begin{proof}
対角な$A$（固有値$\{a_j\}$）で十分。$\Gamma(A)$は$\Lambda^n\mathcal{H}$上で
固有値$\{a_{j_1}\cdots a_{j_n}:j_1<\cdots<j_n\}$を持つ。したがって
\begin{align}
\Tr_{\mathcal{F}_-}(z^N\Gamma(A))
&= \sum_{n=0}^{\infty}z^n\sum_{j_1<\cdots<j_n}a_{j_1}\cdots a_{j_n}
= \sum_{n=0}^{\infty}z^n\,e_n(a_1,a_2,\ldots)
= \prod_j(1+za_j)\\
&= \det\nolimits_F(1+zA).
\end{align}
ここで$e_n$は$n$次基本対称関数であり、
基本対称関数の積公式$\sum_n e_n(a)z^n=\prod_j(1+za_j)$を用いた。
$A$がトレースクラスであるから積は絶対収束する。
\end{proof}

\begin{lemma}[一粒子射影]\label{lem:projector}
Fock空間における一粒子セクター$\Lambda^1\mathcal{H}\cong\mathcal{H}$
への射影は
\begin{equation}\label{eq:projector}
P_1 = \int_0^1 e^{2\pi iu(N-1)}\,du.
\end{equation}
\end{lemma}

\begin{proof}
$N$の固有値$n\in\Z_{\geq 0}$を持つ任意の固有ベクトルに対して：
$\int_0^1 e^{2\pi iu(n-1)}\,du = \delta_{n,1}$。
これは$\Z$の指標の標準的直交性
$\int_0^1 e^{2\pi imu}\,du=\delta_{m,0}$を$m=n-1$に適用したものである。
\end{proof}

\begin{theorem}[普遍射影行列式定理]\label{thm:universal}
$(\mathcal{H},H)$を自己共役ハミルトニアン$H$を持つ量子系とし、
ある$\beta>0$で$e^{-\beta H}$がトレースクラスであるとする
（全ての有限体積格子多体系はこの条件を満たす）。
$\beta=\beta_X(\tau)=2\pi\tau_2$と置く。

ホロノミー$u\in[0,1)$を持つLoNalogy行列式ファイバーを
\begin{equation}\label{eq:detbundle}
Z_{\det}^X(\tau,u;H)
:= \det\nolimits_F\!\left(1+e^{2\pi iu}\,e^{-\beta_X(\tau)H}\right)
\end{equation}
と定義する。このときギブス分配関数は第1フーリエ・ホロノミーモードである：
\begin{equation}\label{eq:universalthm}
\boxed{Z_H(\tau)
:= \Tr_{\mathcal{H}}\,e^{-\beta_X(\tau)H}
= \int_0^1 e^{-2\pi iu}\,Z_{\det}^X(\tau,u;H)\,du.}
\end{equation}
\end{theorem}

\begin{proof}
$A:=e^{-\beta_X H}$（仮定によりトレースクラス）、$z:=e^{2\pi iu}$と置く。

\emph{ステップ1：一粒子トレースのFock空間表現。}
$\Lambda^1\mathcal{H}\cong\mathcal{H}$かつ$\Lambda^1 A=A$であるから
\[
\Tr_{\mathcal{H}}(e^{-\beta_X H})
= \Tr_{\mathcal{F}_-(\mathcal{H})}(P_1\,\Gamma(A)).
\]

\emph{ステップ2：$P_1$のフーリエ表現の代入。}
補題\ref{lem:projector}により
\[
\Tr_{\mathcal{F}_-}(P_1\,\Gamma(A))
= \int_0^1 e^{-2\pi iu}\,
\Tr_{\mathcal{F}_-}(e^{2\pi iuN}\,\Gamma(A))\,du.
\]
トレースと積分の交換は優収束定理で正当化される。

\emph{ステップ3：Fock空間トレース恒等式の適用。}
補題\ref{lem:focktrace}より
$\Tr_{\mathcal{F}_-}(e^{2\pi iuN}\Gamma(A))
= \det_F(1+e^{2\pi iu}e^{-\beta_X H})$。

\emph{ステップ4：結合。}
$\Tr_{\mathcal{H}}e^{-\beta_X H}
= \int_0^1 e^{-2\pi iu}\det_F(1+e^{2\pi iu}e^{-\beta_X H})\,du$。
\end{proof}

\begin{corollary}[生成汎関数]\label{cor:generating}
任意の自己共役観測量$O$に対して、
$Z_H(J;\tau) := \Tr_{\mathcal{H}}e^{-\beta_X(\tau)(H-JO)}$と置くと
\begin{equation}
Z_H(J;\tau)
= \int_0^1 e^{-2\pi iu}\,
\det\nolimits_F\!\left(1+e^{2\pi iu}\,e^{-\beta_X(\tau)(H-JO)}\right)du
\end{equation}
であり、全ての熱相関関数は
$\langle O^n\rangle_\tau
= Z_H(0;\tau)^{-1}(\partial^n/\partial J^n)Z_H(J;\tau)|_{J=0}$
で生成される。
\end{corollary}

\begin{proof}
定理\ref{thm:universal}において$H$を$H-JO$に置き換える。
\end{proof}


\section{単一ファイバーでは不十分}\label{sec:singlefiber}

\begin{proposition}[単一ファイバーの不十分性]\label{prop:singlefiber}
一般に$\Tr_{\mathcal{H}}e^{-\beta H}
\neq \det_F(1+ze^{-\beta H})$（固定の$z$に対して）。
\end{proposition}

\begin{proof}
$H$の固有値を$\{E_j\}$とすると
$\det_F(1+ze^{-\beta H})
= \prod_j(1+ze^{-\beta E_j})
= 1 + z\sum_j e^{-\beta E_j} + z^2\sum_{i<j}e^{-\beta(E_i+E_j)} + \cdots$。
ギブス分配関数は$z^1$の係数であり、行列式全体ではない。
高次項は一般に非零。
\end{proof}


%% ============================================================
\part{低ランク単一ファイバー圧縮のNo-Go定理}
%% ============================================================

\section{主張と証明}\label{sec:nogo}

\begin{theorem}[低ランク圧縮のno-go]\label{thm:nogo}
$(\mathcal{H},H)$を$\dim\mathcal{H}=D$のgenericな非縮退スペクトル
$\{E_1<E_2<\cdots<E_D\}$を持つ量子系とする。
$\dim\mathcal{K}=r$の補助ヒルベルト空間と$\mathcal{K}$上の
自己共役作用素$K$が存在して
\begin{equation}
\Tr_{\mathcal{H}}e^{-\beta H}
= \int_0^1 e^{-2\pi iu}\,
\det_F\!\left(1+e^{2\pi iu}e^{-\beta K}\right)du
\end{equation}
が全ての$\beta>0$で成り立つならば、$r\geq D$。
Generic非縮退$H$では等号$r=D$が必要であり、$K$の固有値は$H$の
それと（多重集合として）一致しなければならない。
\end{theorem}

\begin{proof}
\emph{ステップ1：第1ホロノミー係数の抽出。}
ホロノミー積分は$\det_F(1+ze^{-\beta K})$の$z^1$係数を
抽出するから、仮定は
$\sum_{a=1}^D e^{-\beta E_a} = \sum_{j=1}^r e^{-\beta\varepsilon_j}$
（全$\beta>0$）に帰着する。

\emph{ステップ2：指数和の一意性。}
両辺は離散測度のラプラス変換：
$\mu_H = \sum_{a=1}^D \delta_{E_a}$、
$\mu_K = \sum_{j=1}^r \delta_{\varepsilon_j}$。
ラプラス変換は$[0,\infty)$上の有限ボレル測度の空間で単射であるから
$\mu_H=\mu_K$。

\emph{ステップ3：ランク下限。}
Generic非縮退$H$では$D$個のdistinctなレベルがあるから$r\geq D$。
\end{proof}

\begin{remark}[解釈]
この定理は射影行列式表現（定理\ref{thm:universal}）が最適であることを
示す：ホロノミー射影はフルランク$D$の補助空間なしには除去できない。
束構造は表現の弱点ではなく本質的特徴である。
\end{remark}


%% ============================================================
\part{転送クラスと局所性保存還元}
%% ============================================================

\section{転送クラス}\label{sec:transfer}

\begin{definition}[転送クラス]\label{def:transfer}
$\Lambda=[1,L]\times B$を伝播方向$L$サイト、横断面$B$の格子とする。
分配関数$Z_\Lambda(\beta)$が\emph{厳密転送クラス}に属するとは、
有限次元の補助（境界）空間$\mathcal{B}_\partial$と線形作用素
$T_\beta:\mathcal{B}_\partial\to\mathcal{B}_\partial$が存在して
$Z_\Lambda(\beta) = \Tr_{\mathcal{B}_\partial}T_\beta^{\,L}$
となることをいう。
\emph{ブリッジランク}は$\chi:=\dim\mathcal{B}_\partial$。
\end{definition}

\begin{proposition}[転送クラス = 局所性保存LoN行列式クラス]\label{prop:transferLoN}
$Z_\Lambda(\beta)$が転送行列$T_\beta$、ブリッジランク$\chi$で
厳密転送クラスに属するならば
\begin{equation}
\boxed{Z_\Lambda(\beta)
= \int_0^1 e^{-2\pi iu}\,
\det\nolimits_F\!\left(1+e^{2\pi iu}\,T_\beta^{\,L}\right)du.}
\end{equation}
\end{proposition}

\begin{proof}
定理\ref{thm:universal}を$\mathcal{H}=\mathcal{B}_\partial$、
$A=T_\beta^L$に適用。
\end{proof}


\section{有限レンジ古典系}\label{sec:classical}

\begin{theorem}[古典局所系は転送クラスに属する]\label{thm:classical}
$d$次元有限格子上の局所状態数$q$、伝播方向の相互作用レンジ$R$、
横断面サイズ$|B|$の有限レンジ古典スピン系に対して、
\begin{equation}
\dim\mathcal{B}_\partial \leq q^{R|B|}
\end{equation}
を満たす厳密転送行列$T_\beta$が存在し、
\begin{equation}
\boxed{Z_\Lambda(\beta) = \int_0^1 e^{-2\pi iu}\,
\det\nolimits_F\!\left(1+e^{2\pi iu}\,T_\beta^{\,L-R+1}\right)du.}
\end{equation}
\end{theorem}

\begin{proof}
ハミルトニアンを伝播方向に沿って分解する。レンジ$R$のため、
各局所エネルギーは高々$R$連続スライスにしか依存しない。
「最後の$R$枚のスライスの配置」を補助状態と見なせば、
1ステップ更新は有限行列$T_\beta$（次元$\leq q^{R|B|}$）で表せる。
周期境界条件下で分配関数は$(L-R+1)$乗のトレースになる。
命題\ref{prop:transferLoN}を適用してLoN行列式表現を得る。

この還元は真正である：体積$q^{L|B|}$の和を、
境界サイズ$q^{R|B|}$の転送問題に落としている。
\end{proof}


\section{例：一次元イジングモデル}\label{sec:ising}

\begin{corollary}[LoN行列式クラスの1Dイジング]\label{cor:ising}
一次元イジングモデル
\[
H(\sigma) = -J\sum_{i=1}^N\sigma_i\sigma_{i+1}
-h\sum_{i=1}^N\sigma_i
\qquad(\sigma_i=\pm1\text{、周期境界})
\]
は転送行列
\begin{equation}
T = \begin{pmatrix}
e^{\beta J+\beta h} & e^{-\beta J}\\
e^{-\beta J} & e^{\beta J-\beta h}
\end{pmatrix}
\end{equation}
を持ち、ブリッジランク$\chi=2$で
\begin{equation}\label{eq:isingLoN}
\boxed{Z_N(\beta)
= \int_0^1 e^{-2\pi iu}\,
\det\!\left(1+e^{2\pi iu}\,T^N\right)du.}
\end{equation}
\end{corollary}

\begin{proof}
$T$の固有値を$\lambda_\pm$とする。明示的に計算すると
\begin{align}
Z_N &= \Tr T^N = \lambda_+^N+\lambda_-^N\\
&= \int_0^1 e^{-2\pi iu}\,
(1+e^{2\pi iu}\lambda_+^N)(1+e^{2\pi iu}\lambda_-^N)\,du.
\end{align}
展開すると
\begin{align}
&= \int_0^1 e^{-2\pi iu}\bigl[1
+e^{2\pi iu}(\lambda_+^N+\lambda_-^N)
+e^{4\pi iu}\lambda_+^N\lambda_-^N\bigr]\,du.
\end{align}
指標の直交性により$e^{-2\pi iu}$と$e^{2\pi iu}$の項は
0に積分され、
$= \lambda_+^N+\lambda_-^N = Z_N$。$\checkmark$
\end{proof}

\begin{remark}
これは「イジングモデルをLoNalogyから厳密に導出できるか？」
への具体的回答である。できる。1Dイジング分配関数はランク2の
LoN行列式束の第1ホロノミーモードである。
\end{remark}


\section{境界ランクの最適性}\label{sec:optimality}

\begin{theorem}[境界ランクの最適性]\label{thm:optimality}
幅$|B|$のストリップ上のgenericな結合定数を持つ
最近接$q$-状態モデルに対して、ブリッジランクは
\begin{equation}
\boxed{\chi \geq q^{|B|}.}
\end{equation}
\end{theorem}

\begin{proof}
最近接モデル（$R=1$）の行--行転送行列$T_\beta$は
$q^{|B|}\times q^{|B|}$行列である。Genericな結合定数に対して
$T_\beta$は$q^{|B|}$個のdistinctな非零固有値$\lambda_1,\ldots,\lambda_{q^{|B|}}$
を持つ。

$\chi<q^{|B|}$次元の補助作用素$A$（固有値$\mu_1,\ldots,\mu_\chi$）が
全$L\geq 1$で$\Tr T_\beta^L=\Tr A^L$を満たすと仮定する。
べき和$p_L(T)=\sum_j\lambda_j^L=\sum_k\mu_k^L=p_L(A)$が全$L$で一致する。
ニュートンの恒等式により、べき和は基本対称多項式を一意に決定し、
したがって特性多項式を決定する。
$T$の非零固有値が$q^{|B|}$個distinctであるから、
$\chi<q^{|B|}$の固有値でマッチすることは不可能。$\chi\geq q^{|B|}$。
\end{proof}


%% ============================================================
\part{量子多体系への拡張}
%% ============================================================

\section{トロッター埋め込み}\label{sec:trotter}

\begin{theorem}[トロッター埋め込みによる量子系]\label{thm:trotter}
$(\mathcal{H},H)$を$H=H_A+H_B$の有限体積量子多体系とする。
正整数$m$（トロッター数）に対して
$Z_m(\beta) := \Tr_{\mathcal{H}}(e^{-\beta H_A/m}e^{-\beta H_B/m})^m$
とすると：
(i) $Z_m(\beta)\to Z(\beta)=\Tr e^{-\beta H}$（$m\to\infty$、トロッター積公式）。
(ii) 各有限$m$で$Z_m(\beta)=\Tr T_{\beta,m}^m$
（$T_{\beta,m}=e^{-\beta H_A/m}e^{-\beta H_B/m}$）。
(iii) したがって各有限$m$で
\begin{equation}
\boxed{Z_m(\beta) = \int_0^1 e^{-2\pi iu}\,
\det\nolimits_F\!\left(1+e^{2\pi iu}\,T_{\beta,m}^{\,m}\right)du.}
\end{equation}
\end{theorem}

\begin{proof}
(i)はトロッター積公式（トレースノルム収束）。
(ii)は$T_{\beta,m}:=e^{-\beta H_A/m}e^{-\beta H_B/m}$の定義から直接。
(iii)は$A=T_{\beta,m}^m$に定理\ref{thm:universal}を適用。
\end{proof}

\begin{corollary}[LoN束の極限としての量子分配関数]\label{cor:quantumlimit}
任意の有限体積量子多体系に対して：
\begin{equation}
\Tr\,e^{-\beta H}
= \lim_{m\to\infty}\int_0^1 e^{-2\pi iu}\,
\det\nolimits_F\!\left(1+e^{2\pi iu}\,T_{\beta,m}^{\,m}\right)du.
\end{equation}
極限と積分の交換は優収束定理で正当化される。
\end{corollary}

\begin{proof}
定理\ref{thm:trotter}(i)と(iii)を結合。
\end{proof}


%% ============================================================
\part{完全分類}
%% ============================================================

\section{結果の要約}\label{sec:summary}

\begin{center}
\begin{tabular}{p{5.0cm}lll}
\toprule
結果 & 参照 & レベル\\
\midrule
\multicolumn{3}{l}{\textbf{ガウス還元核（第I部）}}\\
\quad シューア = ガウス消去 & 定理\ref{thm:KT} & A\\
\quad 三セクターシューア & 定理\ref{thm:threesector} & A\\
\quad 統計力学は特殊例 & 命題\ref{prop:statmech} & A\\
\midrule
\multicolumn{3}{l}{\textbf{ボルツマン・松原・KMS・ギブス（第II--V部）}}\\
\quad ボルツマン（$q$-展開） & 定理\ref{thm:boltzmann} & A\\
\quad 松原（スピン構造） & 定理\ref{thm:matsubara} & A\\
\quad KMS（戸田--竹崎） & 定理\ref{thm:kms} & A\\
\quad ギブス状態 & 定理\ref{thm:gibbs} & A\\
\quad ギブス変分原理 & 定理\ref{thm:gibbsvar} & A\\
\quad ホーキング温度 & 命題\ref{prop:hawking} & A条件付\\
\midrule
\multicolumn{3}{l}{\textbf{普遍埋め込み（第VI部）}}\\
\quad 射影行列式の普遍性 & 定理\ref{thm:universal} & A\\
\quad 生成汎関数 & 系\ref{cor:generating} & A\\
\quad 単一ファイバーの不十分性 & 命題\ref{prop:singlefiber} & A\\
\midrule
\multicolumn{3}{l}{\textbf{No-goと最適性（第VII--VIII部）}}\\
\quad 低ランク圧縮のno-go & 定理\ref{thm:nogo} & A\\
\quad 古典局所系$\Rightarrow$転送クラス & 定理\ref{thm:classical} & A\\
\quad 1Dイジング & 系\ref{cor:ising} & A\\
\quad 境界ランク最適性 & 定理\ref{thm:optimality} & A\\
\midrule
\multicolumn{3}{l}{\textbf{量子拡張（第IX部）}}\\
\quad トロッター埋め込み & 定理\ref{thm:trotter} & A\\
\quad LoN束の量子極限 & 系\ref{cor:quantumlimit} & A\\
\bottomrule
\end{tabular}
\end{center}

\section{4部分類}\label{sec:fourpart}

LoNalogy--統計力学の普遍性に関する完全な判定は：

\begin{enumerate}[label=\textbf{(\Alph*)}]
\item \textbf{厳密スペクトル普遍性：真。}
全ての有限体積ギブス多体系はLoN行列式束の厳密な第1ホロノミーモード
（定理\ref{thm:universal}）。

\item \textbf{普遍的低ランク単一ファイバー圧縮：偽。}
$r<D$の補助系ではgenericな$D$次元分配関数を全$\beta$で
厳密再現不可能（定理\ref{thm:nogo}）。

\item \textbf{厳密局所普遍性：真。}
有限レンジ局所古典系は厳密に転送クラスであり、
従って局所性保存LoN行列式クラス（定理\ref{thm:classical}）。

\item \textbf{境界ランク最適性：真。}
Generic局所モデルは境界ランクスケーリング$\chi\geq q^{|B|}$を
要求し、これは普遍的には改善不可能（定理\ref{thm:optimality}）。
\end{enumerate}

一つの公式で：
\begin{equation}
\boxed{\parbox{0.85\textwidth}{\centering
LoN行列式普遍性は射影形式で真、\\
低ランク単一ファイバーで偽、\\
転送境界形式で厳密かつ最適。}}
\end{equation}


%% ============================================================
\section*{結論}
\addcontentsline{toc}{section}{結論}
%% ============================================================

LoNalogy行列式セクターと統計力学の完全な関係を確立した。
結果は二つの分岐を持つ論理鎖を形成する。

\textbf{分岐1：なぜギブスか？}
\begin{equation}
X(2)\;\to\; Q=e^{-2\pi\tau_2}\;\to\; e^{-\beta E}
\;\to\;\text{松原}\;\to\;\text{KMS}\;\to\;\text{ギブス}。
\end{equation}
ボルツマン因子は上半平面のノームである。ギブス状態は
$X(2)$-モジュラーKMS状態の一意なエントロピー最大化子である。
いずれも公準ではない。

\textbf{分岐2：普遍性はどこまで行くか？}
\begin{equation}
\text{射影束：厳密}
\;\xrightarrow{\;\text{no-go}\;}
\text{単一ファイバー：フルランク必要}
\;\xrightarrow{\;\text{転送}\;}
\text{局所系：最適}。
\end{equation}
全てのギブス系はLoN行列式束に埋め込まれる。
普遍的低ランク圧縮は存在しない。
局所系は転送クラスを通じて最適な境界ランクを達成する。
量子系はトロッター埋め込みにより含まれる。

4部分類(A)--(D)は完全である：普遍性の問題にはopenな部分はない。
$X(2)$から平衡統計力学の全構造への論理鎖は閉じている。

一文で：

\begin{quote}
\emph{ギブスは公理ではない。$X(2)$のモジュラー定理であり、
ボルツマンの指数関数は上半平面のノームである。
全てのギブス系は射影行列式束の中に住み、
局所系は転送クラスで最適な境界ランクを達成する。}
\end{quote}


%% ============================================================
\section*{謝辞}
\addcontentsline{toc}{section}{謝辞}
%% ============================================================

著者はClaude (Anthropic)に、導出の検証、証明の組織化、原稿作成への
支援を感謝する。ボルツマン因子が$X(2)$の$q$-一様化であるという
重要な観察は、LoNalogy行列式セクターの共同分析を通じて
同定された。No-go定理と転送クラス分類は同一の共同セッションで
開発された。


%% ============================================================
\begin{thebibliography}{99}

\bibitem{LoNalogyv87}
M.~Nakamura,
\textit{LoNalogy v87.1: Elliptic Modular Framework for Quantum Field Theory
and Millennium Problems},
Zenodo (2026), \href{https://doi.org/10.5281/zenodo.19115338}{doi:10.5281/zenodo.19115338}.

\bibitem{YMmassGap}
M.~Nakamura,
\textit{Yang--Mills Existence and Mass Gap: A Constructive Proof via
Elliptic Modular Geometry and Transfer-Poincar\'e Descent},
Zenodo (2026), \href{https://doi.org/10.5281/zenodo.19183838}{doi:10.5281/zenodo.19183838}.

\bibitem{ParticleCosmo}
M.~Nakamura,
\textit{Standard Model Structure, GUT Scale, and Dark Sector from Elliptic
Modular Geometry: Predictions of LoNalogy},
姉妹論文 (2026).

\bibitem{BratteliRobinson}
O.~Bratteli and D.~W.~Robinson,
\textit{Operator Algebras and Quantum Statistical Mechanics},
Vol.~2, 2nd ed., Springer (1997).

\bibitem{Simon}
B.~Simon,
\textit{Trace Ideals and Their Applications},
2nd ed., Mathematical Surveys and Monographs \textbf{120},
AMS (2005).

\bibitem{Haag}
R.~Haag,
\textit{Local Quantum Physics: Fields, Particles, Algebras},
2nd ed., Springer (1996).

\bibitem{Boltzmann1877}
L.~Boltzmann,
\textit{\"Uber die Beziehung zwischen dem zweiten Hauptsatze der
mechanischen W\"armetheorie und der Wahrscheinlichkeitsrechnung
respektive den S\"atzen \"uber das W\"armegleichgewicht},
Wien.\ Ber.\ \textbf{76}, 373--435 (1877).

\bibitem{Jaynes1957}
E.~T.~Jaynes,
\textit{Information Theory and Statistical Mechanics},
Phys.\ Rev.\ \textbf{106}, 620--630 (1957).

\bibitem{Baxter}
R.~J.~Baxter,
\textit{Exactly Solved Models in Statistical Mechanics},
Academic Press (1982).

\bibitem{Trotter1959}
H.~F.~Trotter,
\textit{On the Product of Semi-Groups of Operators},
Proc.\ Amer.\ Math.\ Soc.\ \textbf{10}, 545--551 (1959).

\end{thebibliography}

\end{document}
